\section{Generative Adversarial Networks (GANs)}

\subsection{De la modelización explícita al muestreo implícito}

El VAE intenta modelar \textbf{explícitamente} la densidad de probabilidad de los datos $p(x)$. Pero si nuestro objetivo es solo \textbf{generar muestras nuevas}, ¿es necesario modelar explícitamente la densidad?

\begin{figure}[H]
\centering
\includegraphics[width=0.85\textwidth]{slides-images/slide-48.jpg}
\caption{La idea de la GAN: sin modelar explícitamente la densidad, aprender directamente una función que mapea el ruido a los datos\protect\footnotemark}
\end{figure}
\footnotetext{Diapositivas, página 48.}

\begin{knowledgebox}{La filosofía de diseño de la GAN}
La idea central de la GAN es: en lugar de modelar explícitamente $p(x)$, aprender directamente una \textbf{transformación} de una distribución simple (como el ruido gaussiano) a una distribución de datos compleja. En términos concretos, se muestrea de una distribución de ruido simple y luego se entrena una red generadora que transforma el ruido en muestras de datos verosímiles. La pregunta clave se convierte en: ¿cómo entrenar este generador?
\end{knowledgebox}

\subsection{La idea del entrenamiento adversario de la GAN}

La GAN propone una solución elegante: introducir dos redes que compiten entre sí: el \textbf{generador} (Generator, $G$) y el \textbf{discriminador} (Discriminator, $D$).

\begin{itemize}
    \item \textbf{Generador $G$}: recibe ruido aleatorio $z$ e intenta generar datos falsos lo más verosímiles posible $x_{\text{fake}} = G(z)$
    \item \textbf{Discriminador $D$}: recibe una muestra de datos (que puede ser real o generada) y produce la probabilidad de que esa muestra sea un dato real $D(x) \in {[}0, 1{]}$
\end{itemize}

\begin{figure}[H]
\centering
\includegraphics[width=0.85\textwidth]{slides-images/slide-68.jpg}
\caption{La arquitectura de entrenamiento de la GAN: el enfrentamiento entre el generador y el discriminador\protect\footnotemark}
\end{figure}
\footnotetext{Diapositivas, página 68.}

Las dos redes forman un \textbf{juego adversario}:
\begin{itemize}
    \item $G$ intenta generar datos falsos capaces de engañar a $D$
    \item $D$ intenta distinguir correctamente los datos reales de los datos falsos generados por $G$
    \item Cuando el entrenamiento alcanza el óptimo global, $G$ es capaz de generar muestras consistentes con la distribución de los datos reales
\end{itemize}

\subsection{La intuición del entrenamiento de la GAN}

El curso mostró de manera vívida el proceso de entrenamiento de la GAN con una animación sobre una línea de datos unidimensional:

\begin{figure}[H]
\centering
\includegraphics[width=0.85\textwidth]{slides-images/slide-50.jpg}
\caption{La intuición del entrenamiento de la GAN: el generador parte del ruido y se acerca gradualmente a la distribución de los datos reales\protect\footnotemark}
\end{figure}
\footnotetext{Diapositivas, página 50.}

\textbf{La iteración del proceso de entrenamiento}:
\begin{enumerate}
    \item Al inicio, el generador produce datos falsos de muy baja calidad a partir de ruido aleatorio
    \item El discriminador aprende a distinguir los datos reales (puntos verdes) de los datos falsos (puntos rosas), asignando una probabilidad alta a los datos reales
    \item El generador ajusta la distribución de los datos que produce de acuerdo con la retroalimentación del discriminador, para que los datos falsos se acerquen más a los datos reales
    \item El discriminador se actualiza de nuevo para adaptarse a los datos falsos, ahora mejores
    \item Se itera repetidamente, hasta que los datos falsos resultan indistinguibles de los datos reales
\end{enumerate}

\begin{figure}[H]
\centering
\includegraphics[width=0.85\textwidth]{slides-images/slide-65.jpg}
\caption{Etapa avanzada del entrenamiento: la distribución de los datos falsos (rosa) se acerca a la de los datos reales (verde)\protect\footnotemark}
\end{figure}
\footnotetext{Diapositivas, página 65.}

\subsection{La función de pérdida de la GAN}

El objetivo de entrenamiento de la GAN puede expresarse como un juego \textbf{minimax}:

$$\min_G \max_D \; \mathbb{E}_{z,x}\left[\log D(x) + \log(1 - D(G(z)))\right]$$

\begin{figure}[H]
\centering
\includegraphics[width=0.85\textwidth]{slides-images/slide-70.jpg}
\caption{La función de pérdida de la GAN: el generador minimiza y el discriminador maximiza el mismo objetivo\protect\footnotemark}
\end{figure}
\footnotetext{Diapositivas, página 70.}

\begin{importantbox}{Interpretación de la función de pérdida de la GAN}
\textbf{El objetivo del discriminador (maximizar)}:
\begin{itemize}
    \item $\log D(x)$: para los datos reales $x$, se desea que $D(x) \to 1$ (identificarlos correctamente como reales)
    \item $\log(1 - D(G(z)))$: para los datos generados $G(z)$, se desea que $D(G(z)) \to 0$ (identificarlos correctamente como falsos)
\end{itemize}

\textbf{El objetivo del generador (minimizar)}:
\begin{itemize}
    \item Se desea que $D(G(z)) \to 1$, es decir, que el discriminador clasifique erróneamente los datos falsos como reales
    \item Lo cual equivale a minimizar $\log(1 - D(G(z)))$
\end{itemize}
\end{importantbox}

\begin{warningbox}{La inestabilidad del entrenamiento de la GAN}
En la práctica, el entrenamiento adversario de la GAN puede ser muy inestable. Las dos redes necesitan mantener un progreso de entrenamiento relativamente equilibrado: si el discriminador es demasiado fuerte, el generador no recibe señales de gradiente útiles; si el generador es demasiado fuerte, el discriminador no puede ofrecer retroalimentación eficaz. Los problemas comunes incluyen el colapso de modos (Mode Collapse) y las oscilaciones durante el entrenamiento. Trabajos posteriores, como Wasserstein GAN (WGAN), propusieron funciones de pérdida alternativas más estables.
\end{warningbox}

\subsection{La GAN como transformador de distribuciones}

Una vez terminado el entrenamiento, se descarta el discriminador y solo se usa el generador para crear datos nuevos:

\begin{figure}[H]
\centering
\includegraphics[width=0.85\textwidth]{slides-images/slide-72.jpg}
\caption{La GAN ya entrenada: usar solo el generador para crear datos nuevos a partir de ruido\protect\footnotemark}
\end{figure}
\footnotetext{Diapositivas, página 72.}

En el fondo, el generador de la GAN aprende una \textbf{transformación de distribuciones}: una función que mapea de la distribución de ruido gaussiano $z \sim \mathcal{N}(0, 1)$ a la distribución de los datos objetivo.

\begin{figure}[H]
\centering
\includegraphics[width=0.85\textwidth]{slides-images/slide-75.jpg}
\caption{La GAN es un transformador de distribuciones: una función que mapea de la distribución de ruido a la distribución de los datos\protect\footnotemark}
\end{figure}
\footnotetext{Diapositivas, página 75. Fuente: Brock+ ICLR 2019.}

Al realizar una interpolación suave en el espacio de ruido, las muestras generadas también hacen una transición suave sobre la variedad de los datos, lo cual muestra que el generador aprendió la estructura continua de los datos.

\subsection{El desarrollo y las aplicaciones de la GAN}

\textbf{Progressive Growing of GANs}: al aumentar gradualmente el número de capas de la red y la resolución, se pueden generar imágenes de rostros de alta resolución extremadamente verosímiles.

\begin{figure}[H]
\centering
\includegraphics[width=0.85\textwidth]{slides-images/slide-78.jpg}
\caption{Rostros de alta resolución generados por Progressive GAN (todos generados por IA)\protect\footnotemark}
\end{figure}
\footnotetext{Diapositivas, página 78. Fuente: Karras+ ICLR 2018.}

\textbf{Conditional GAN y Pix2Pix}: al incorporar información condicional en el generador y el discriminador, se pueden realizar tareas de traducción de imágenes pareadas, como generar fotografías realistas de escenas urbanas a partir de mapas de segmentación semántica.

\begin{figure}[H]
\centering
\includegraphics[width=0.85\textwidth]{slides-images/slide-80.jpg}
\caption{Conditional GAN (Pix2Pix): traducción de imágenes pareadas\protect\footnotemark}
\end{figure}
\footnotetext{Diapositivas, página 80. Fuente: Isola+ CVPR 2017.}

\begin{figure}[H]
\centering
\includegraphics[width=0.85\textwidth]{slides-images/slide-82.jpg}
\caption{Resultados de aplicación de Pix2Pix: conversión bidireccional entre mapas y fotografías aéreas\protect\footnotemark}
\end{figure}
\footnotetext{Diapositivas, página 82. Fuente: Isola+ CVPR 2017.}

\textbf{CycleGAN}: permite la conversión entre dominios sin necesidad de datos pareados. No se limita a las imágenes; también puede aplicarse a ámbitos como la conversión de voz.

\begin{figure}[H]
\centering
\includegraphics[width=0.85\textwidth]{slides-images/slide-85.jpg}
\caption{La aplicación de CycleGAN en la conversión de voz\protect\footnotemark}
\end{figure}
\footnotetext{Diapositivas, página 85.}

\begin{knowledgebox}{La innovación de CycleGAN}
La GAN tradicional genera la distribución objetivo a partir de una distribución de ruido. CycleGAN amplía esta idea: en lugar de partir del ruido, parte de una distribución de datos y aprende una función que mapea a otra distribución de datos. La innovación clave es el uso de la \textbf{pérdida de consistencia cíclica} (Cycle Consistency Loss), que garantiza que, al convertir del dominio A al dominio B y de vuelta al dominio A, se recupere la entrada original. Esto permite entrenar sin necesidad de datos pareados.
\end{knowledgebox}

\subsection{Resumen de la sección}

Mediante el entrenamiento adversario entre el generador y el discriminador, la GAN logra el objetivo de generar muestras de alta calidad sin necesidad de modelar explícitamente la densidad. En esencia, la GAN es un transformador de distribuciones que mapea una distribución de ruido simple a una distribución de datos compleja. Aunque el entrenamiento de la GAN presenta el reto de la inestabilidad, su calidad de generación (especialmente en imágenes) suele superar a la del VAE, y de ella se han derivado variantes poderosas como Conditional GAN y CycleGAN.


